\documentclass[10pt,a4paper]{amsart}
\usepackage[margin=19mm]{geometry}
\usepackage{amsmath,amssymb,mathtools}
\usepackage[scr=rsfs,cal=euler]{mathalfa}
\usepackage{xfrac}
\usepackage[unicode,hidelinks,bookmarks=false]{hyperref}
\newcommand{\ie}{i.e.\ }
\newcommand{\cf}{cf.\ }
\newcommand{\resp}{respectively\ }
\newcommand{\Subsection}{\subsection}
\providecommand{\nobreakdash}{\nobreak\hbox{-}}
\setlength{\emergencystretch}{2em}
\multlinegap0pt
\usepackage{tikz}
\usepackage{enumerate}
\usepackage{mathtools}
\usepackage{bm}
\newtheorem{proposition}{Proposition}
\newtheorem{lemma}{Lemma}
\newcommand{\Tr}{\mathrm{Tr}}
\newcommand{\cA}{\mathcal{A}}
\newcommand{\cL}{\mathcal{L}}
\newcommand{\cP}{\mathcal{P}}
\newcommand{\cov}{\mathrm{Cov}}
\newcommand{\rE} {\mathbb{E}}
\newcommand{\sN}{{\scriptscriptstyle N}}
\newcommand{\tr}{\mathrm{tr}}
\title{infinitesimal freeness for orthogonally\\[5pt] invariant random matrices}
\author{Guillaume C\'ebron (Toulouse), James A Mingo (Queen's)}
\date{Source: arXiv:2506.05139v2 (CC BY 4.0)}
\begin{document}
\maketitle

\noindent\emph{Source and licence.} infinitesimal freeness for orthogonally\\[5pt] invariant random matrices. Version arXiv:2506.05139v2 (CC BY 4.0), distributed by arXiv under the Creative Commons Attribution 4.0 International licence, http://creativecommons.org/licenses/by/4.0/, as recorded in the arXiv licence metadata for this exact version and shown on the abstract page https://arxiv.org/abs/2506.05139v2. Full licence notice: \texttt{licenses/arxiv-2506.05139-CC-BY-4.0.txt}.\par\smallskip

\noindent\textit{Editorial scope.} Counted unit: the article's lemma induction (source0.tex lines 1488--1504). The article's complete original proof of it is delivered with it (source0.tex lines 1506--1665, one \texttt{proof} environment of 160 lines). No mathematics is written, repaired or re-derived: the statement and the proof are the article's own lines, in the article's own notation, and no retained line is substituted or re-flowed.  Retained uncounted as the source-local context the proof needs: lines 1051--1074; lines 1241--1243; lines 1245--1247; lines 1249--1251; lines 1300--1312; lines 1473--1484.  Nothing was omitted from any retained span, so the package carries no omission record.  No retained line contains a citation, so the wrapper carries no bibliography and no bibliography entry.  No retained line contains a figure environment, an image-inclusion command or drawing code, so no image is delivered and the package declares no asset dependency.  Editorial formatting only, in addition: an AMS \texttt{amsart} preamble that restates the article's own declarations and macros used by the retained lines, namely the environments \texttt{proposition}, \texttt{lemma} on separate counters, together with the standard packages those lines need.  The retained environments print in this document as Proposition 1, Lemma 1 in order of appearance, whereas the article prints the same items with the numbering of its own full document.  The complete native source of the article as distributed in the arXiv e-print for this exact version is bundled unchanged as \texttt{sources/179384/source0.tex}.
\begin{proposition}\label{prop:int_by_parts}
Let $O_N$ be a Haar distributed orthogonal random matrix of
size $N$, $n$ even, and $M_1,\ldots,M_n\in
M_N(\mathbb{C})$. We have
\begin{align*}\lefteqn{
(N-1)\cdot \mathbb{E}\big[\Tr(O_NM_1O_N^t\cdot  M_2 \cdot O_NM_3O_N^t\cdots O_NM_{n-1}O_N^t\cdot M_n)\big]  }\\
=&- \kern-0.5em
\mathop{\sum_{k = 1}^{n-1}}_{k\mathrm{\ odd}}
%
\mathbb{E}\big[\Tr(O_N M_1 O_N^t \cdots O_N M_{k} O_N^t \cdot (M_{k+1}\cdots  O_NM_{n-1}O_N^t\cdot M_n)^t)\big]\\
+&
\mathop{\sum_{k = 1}^{n-1}}_{k\mathrm{\ odd}}
%
\mathbb{E}\big[\Tr(O_N M_1O_N^t\cdots O_NM_{k}O_N^t)\cdot \Tr(M_{k+1}\cdots  O_NM_{n-1}O_N^t\cdot M_n)\big]\\
+&
\mathop{\sum_{k = 3}^{n-1}}_{k\mathrm{\ odd}}
%
\mathbb{E}\big[\Tr(O_N M_1O_N^t\cdots M_{k-1}\cdot (O_NM_{k}O_N^t\cdots O_NM_{n-1}O_N^t\cdot  M_n)^t)\big]\\
-&
\mathop{\sum_{k = 3}^{n-1}}_{k\mathrm{\ odd}}
%
\mathbb{E}\big[\Tr(O_N M_1O_N^t\cdots M_{k-1})\cdot \Tr(O_NM_{k}O_N^t\cdots O_NM_{n-1}O_N^t\cdot  M_n)\big]
\end{align*}
\end{proposition}

\begin{equation}\label{eq:assumption one}
\rE(\tr(P_i)) = \tau(p_i) + N^{-1} \tau'(p_i) + o(N^{-1}),
\end{equation}

\begin{equation}\label{eq:assumption two}
k_2(\Tr(P_1), \Tr(P_2)) \to \tau_2(p_1, p_2) \mbox{\ as\ } N \to \infty
\end{equation}

\begin{equation}\label{eq:assumption three}
k_r(\Tr(P_1), \dots, \Tr(P_r)) = o(1) \quad\mbox{for\ } r \geq 3
\end{equation}

\begin{lemma}\label{lemma:outside induction}
Let $\cA_{1, \sN}, \dots, \cA_{s, \sN} \subseteq
M_N(\cL^{\infty-})$ be unital subalgebras such that the
entries of matrices from different subalgebras form
independent sets. Suppose that all, or all but one, of the
subalgebras are orthogonally invariant, and suppose that
each of the subalgebras satisfies $(\ref{eq:assumption
  one})$, $(\ref{eq:assumption two})$ and
$(\ref{eq:assumption three})$. Then the subalgebra generated
by $\cA_{1, \sN}, \dots, \cA_{s, \sN}$ satisfies
$(\ref{eq:assumption one})$, $(\ref{eq:assumption two})$ and
$(\ref{eq:assumption three})$.
\end{lemma}

\begin{lemma}\label{lemma:first inside induction}
Let $P_1, \dots, P_r \in M_N(\cL^{\infty-})$ be such that
$P_i \in \cA_{j_i}$ with $j_1 \not= j_2$, $j_2 \not = j_3$,
\dots, $j_{r-1} \not = j_r$, and for each $i$,
$\rE(\Tr(P_i)) = \tau'(p_i) + o(1)$.

Then for $r \geq 1$
\begin{equation}\label{eq:first intermediate limit distribution}
\rE(\Tr(P_1 \cdots P_r))
= O(1)
\end{equation} 
\end{lemma}

\begin{lemma}\label{lemma:inside induction}
Let $P_1, \dots, P_r \in M_N(\cL^{\infty-})$ be such that
$P_i \in \cA_{j_i}$ with $j_1 \not= j_2$, $j_2 \not = j_3$,
\dots, $j_{r-1} \not = j_r$, $j_r \not = j_1$, and for each
$i$, $\rE(\Tr(P_i)) = \tau'(p_i) + o(1)$. Let $V \in
\ker(j)$, with $V = \{l_1, \dots, l_n\}$ be the block
containing $1$ with $1 = l_1$, $l_{m-1} + 1 < l_m$, and $l_n
< r$.

Then for $r > 1$
\begin{equation}\label{eq:intermediate limit distribution}
\rE(\Tr(P_1 \cdots P_r))
=
\mathop{\sum_{m=2}^{n}}
\tau(p_1 \cdots p_{l_m-1} p_r^t \cdots p_{l_m}^t )  + O(N^{-1}).
\end{equation}
\end{lemma}

\begin{proof} 
By traciality we may assume that $\cA_{j_1}$ is orthogonally invariant. 


As in the hypothesis, $ V \in \ker(i) \in \cP(r)$ is the
block containing $1$ and we write $V = \{ l_1, \dots, l_n\}$
with $1 = l_1$, $l_{k-1} + 1 < l_k$ and $l_n < r$ (because
we have assumed that $j_r \not = j_1$). Let
\[
M_1 = P_1, \quad M_2 = P_2 \cdots P_{l_2 -1},\quad M_3 = P_{l_2},\]
and in general 
\[
M_{2k-1} = P_{l_k} \mbox{\ and\ } M_{2k} = P_{l_k + 1} \cdots P_{l_{k+1} - 1}. 
\]
Then 
\[
P_1 \cdots P_r = M_1 M_2 \cdots M_{2n-1} M_{2n}
\]
with $M_1, M_3, \dots, M_{2k-1}$ all in $\cA_{j_1}$ and
$M_2, M_4, \dots, M_{2n}$ all in $\cA_{j_2} \cup \cdots \cup
\cA_{j_r}$. Since $\cA_{j_1}$ is orthogonally invariant we
have
\[
\rE(\Tr(M_1 M_2 \cdots M_{2n-1} M_{2n}) 
=
\rE(\Tr(OM_1 O^t M_2 \cdots O M_{2n-1} O^t M_{2n})).
\]
In the expression on the right each $M_{2k-1}$ has been
replaced by $OM_{2k-1}O^t$ and each $M_{2k}$ has been left
unchanged. In Proposition \ref{prop:int_by_parts} we assumed
that the $M_k$ matrices were constant and here they are
random but independent of the $O$ matrices, so by using
Fubini's theorem we may apply Proposition
\ref{prop:int_by_parts} to obtain that
\begin{align}\lefteqn{
\rE(\Tr(P_1 \cdots P_r)) = \rE(\Tr(M_1 M_2 \cdots M_{2n - 1} M_{2 n}))} \notag \\
&\mbox{} =
\rE(\Tr(OM_1O^t M_2 \cdots OM_{2n - 1}O^t M_{2 n})) \notag \\
& \mbox{} =
\label{eq:first case}
\frac{-1}{N - 1} \mathop{\sum_{k=1}^{2n - 1}}_{k \mathrm{\ odd}}
\rE( \Tr(O M_1 O^t M_2 \cdots O M_kO^t (M_{k + 1} O \cdots O^t M_{2n})^t )) \\
& \mbox{} +
\label{eq:second case}
\frac{1}{N - 1} \mathop{\sum_{k=3}^{2n - 1}}_{k \mathrm{\ odd}}
\rE( \Tr(O M_1 O^t M_2 \cdots O^t M_{k-1} (O M_{k} O^t \cdots O^t M_{2n})^t )) \\
& \mbox{} +
\label{eq:third case}
\frac{1}{N - 1} \mathop{\sum_{k=1}^{2n - 1}}_{k \mathrm{\ odd}}
\rE( \Tr(O M_1 O^t M_2 \cdots O M_{k}O^t) \Tr(M_{k+1} O \cdots O^t M_{2n}) )) \\
& \mbox{} -
\label{eq:fourth case}
\frac{1}{N - 1} \mathop{\sum_{k=3}^{2n - 1}}_{k \mathrm{\ odd}}
\rE( \Tr(O M_1 O^t M_2 \cdots O^t M_{k-1}) \Tr(OM_{k} O^t \cdots O^t M_{2n}) )).
\end{align}

Now let us consider the limit as $N \to \infty$ of each of
the four terms. Let us start with (\ref{eq:first case}). As
$k$ is odd we have $M_k \in \cA_{j_1}$ and $M_{2n} \in
\cA_{j_r}$ with $j_1 \not = j_r$, thus
\begin{align*} \lefteqn{
\frac{1}{N - 1}\rE( \Tr(O M_1 O^t M_2 \cdots O M_kO^t (M_{k + 1} O \cdots O^t M_{2n})^t ))} \\
& \mbox{} =
\frac{1}{N - 1}\rE( \Tr(O M_1 O^t M_2 \cdots O M_kO^t M_{2n}^t O \cdots O^t M_{k+1})) \\
& \mbox{} =
\frac{1}{N - 1}\rE( \Tr(M_1 M_2 \cdots M_k M_{2n}^t  \cdots  M_{k+1}^t)) \\
& \mbox{} =
\frac{1}{N - 1}\rE( \Tr(P_1 P_2 \cdots P_{l_k} P_{r}^t  \cdots  P_{l_k+1}^t)) = O(N^{-1}),\\
\end{align*}
where the last equality holds by Lemma \ref{lemma:outside induction}.




Next consider (\ref{eq:second case}), with $k = 2m - 1 \geq 3$ and $2 \leq m \leq n$.
\begin{align*} \lefteqn{
\frac{1}{N - 1}\rE( \Tr(O M_1 O^t M_2 \cdots O^t M_{k-1} (O M_{k} O^t \cdots O^t M_{2n})^t ))} \\
& \mbox{} =
\frac{1}{N - 1}\rE( \Tr(O M_1 O^t M_2 \cdots O^t M_{2(m-1)} M_{2n}^t O \cdots O M_{2m-1}^t O^t)) \\
& \mbox{} =
\frac{1}{N - 1}\rE( \Tr(M_1 M_2 \cdots M_{2(m-1)} M_{2n}^t  \cdots  M_{2m-1}^t)) \\
& \mbox{} =
\tau(p_1 \cdots p_{l_m-1} p_r^t \cdots p_{l_m}^t ) + O(N^{-1}),
\end{align*}
where we obtain the last equality because, by Lemma
\ref{lemma:outside induction}, Eq.~(\ref{eq:assumption one})
holds:
\begin{align*}\lefteqn{
\rE(\Tr(P_1 \cdots P_{l_m-1} P_r^t \cdots P_{l_m}^t ) } \\
& =
N \tau(p_1 \cdots p_{l_m-1} p_r^t \cdots p_{l_m}^t )
+ 
\tau'(p_1 \cdots p_{l_m-1} p_r^t \cdots p_{l_m}^t )
+
o(1).
\end{align*}
Next consider (\ref{eq:third case}), with $k = 2m - 1 \geq
3$ and $1 \leq m \leq n$.  According to our notation we have
\[
\rE(\Tr(M_1 M_2 \cdots M_{k}))\ \rE(\Tr(M_{k+1}  \cdots  M_{2n}) ))
\]
\[
=
\rE(\Tr( P_1 \cdots P_{l_m}))\ \rE(\Tr( P_{l_m +1} \cdots P_{r})),
\]
and by induction (on $n$) both of these factors are bounded
functions of $N$ by Lemma \ref{lemma:first inside
  induction}. Hence
\begin{align*}\lefteqn{
\frac{1}{N - 1}
\rE( \Tr(O M_1 O^t M_2 \cdots O M_{k}O^t) \Tr(M_{k+1} O \cdots O^t M_{2n}) ))} \\
& =
\frac{1}{N - 1} \cov(\Tr(O M_1 O^t M_2 \cdots O M_{k}O^t), 
                     \Tr(M_{k+1} O \cdots O^t M_{2n}) )) \\
& \quad+
\frac{1}{N - 1} \rE(\Tr(O M_1 O^t M_2 \cdots O M_{k}O^t)) \
                 \rE(\Tr(M_{k+1} O \cdots O^t M_{2n}) )) \\
& =
\frac{1}{N - 1} \rE(\Tr(M_1 M_2 \cdots M_{k})) \
                 \rE(\Tr(M_{k+1}  \cdots  M_{2n}) )) + O(N^{-1}) \\
&=
                 \frac{1}{N - 1} \rE(\Tr(P_1 \cdots P_{l_m})) \
                 \rE(\Tr(P_{l_m+1}  \cdots  P_{r}) )) + O(N^{-1}) \\
& = O(N^{-1}).
\end{align*}
Finally consider (\ref{eq:fourth case}), $k = 2m - 1$ odd
with $m \geq 2$ we have
\[
\rE(\Tr(M_1 M_2 \cdots M_{k-1}))\ \rE(\Tr(M_{k}  \cdots  M_{2n}) ))
\]
\[
=
\rE(\Tr( P_1 \cdots P_{l_m -1 }))\ \rE(\Tr( P_{l_m} \cdots P_{r})).
\]
Now again by Lemma \ref{lemma:first inside induction} , both
of these factors are bounded functions of $N$, hence
\begin{align*}\lefteqn{
\frac{1}{N - 1}
\rE( \Tr(O M_1 O^t M_2 \cdots O^t M_{k-1}) \Tr(O M_{k} O^t \cdots O^t M_{2n}) ))} \\
& =
\frac{1}{N - 1} \cov(\Tr(O M_1 O^t M_2 \cdots O^t M_{k-1}), 
                     \Tr(OM_{k} O^t \cdots O^t M_{2n}) )) \\
& \quad+
\frac{1}{N - 1} \rE(\Tr(O M_1 O^t M_2 \cdots O^t M_{k-1})) \
                 \rE(\Tr(O M_{k} O^t \cdots O^t M_{2n}) )) \\
& =
\frac{1}{N - 1} \rE(\Tr(M_1 M_2 \cdots M_{k-1})) \
                 \rE(\Tr(M_{k}  \cdots  M_{2n}) )) + O(N^{-1}) \\
&=
                 \frac{1}{N - 1} \rE(\Tr(P_1 \cdots P_{l_m-1})) \
                 \rE(\Tr(P_{l_m}  \cdots  P_{r}) )) + O(N^{-1}) \\
& = O(N^{-1}).
\end{align*}
Thus we have 
\[
\rE(\Tr(P_1 \cdots P_r))
=
\sum_{m=2}^n \tau(p_1 \cdots p_{l_m-1} p_r^t \cdots p_{l_m}^t ) + O(N^{-1})
\]
\end{proof} 

\end{document}
